%! Parametric Functions — Unit 8, Lesson 8.1 — Lesson Plan
\documentclass[10pt]{article}
\usepackage{precalculus-article}
\usepackage{precalculus-boxes}
\usepackage{graphicx}
\graphicspath{{images/}}
\newcommand{\TallMath}[1]{$\displaystyle #1\rule[-1.4em]{0pt}{3.2em}$}
\newcommand{\CourseName}{Precalculus}
\newcommand{\MeetingLength}{60 minutes}
\newcommand{\UnitNumberName}{Unit 8: Parametric Functions, Vectors, and Matrices \quad}
\newcommand{\LessonNumberName}{Lesson 8.1: Parametric Functions}

\begin{document}

%----------------------------------------------------------------------
% Title Block
%----------------------------------------------------------------------
\begin{center}
  {\Large\bfseries \CourseName} \\[4pt]
  {\small\bfseries \UnitNumberName \LessonNumberName \quad (2 instructional periods, \MeetingLength\ each)}
\end{center}

%----------------------------------------------------------------------
% Primary Objective
%----------------------------------------------------------------------
\begin{tcolorbox}[colback=sky, colframe=navy, boxrule=0.9pt, arc=2mm,
    left=2mm, right=2mm, top=1.5mm, bottom=1.5mm]
  \textbf{Primary Objective:} Students will construct tables of values and
  graphs for parametric functions by evaluating $x(t)$ and $y(t)$ at selected
  values of the parameter $t$, connecting points in order of increasing $t$,
  and interpreting domain restrictions as start and end points on the curve.
  \textit{(Big Idea: Functions)}
\end{tcolorbox}

\vspace{0.08in}

%----------------------------------------------------------------------
% Priority Ideas & Skills
%----------------------------------------------------------------------
\begin{skillbox}[Priority Ideas \& Skills]{goldbox}
\begin{tabularx}{\linewidth}{@{} p{0.46\linewidth} X @{}}
\textbf{AP Skills} & \textbf{Key Understandings (EKs)} \\
\midrule
\textbf{1.A} \textit{Solve equations analytically} ---
  Solve $x(t) = c$ or $y(t) = c$ for the parameter $t$ to find when the
  curve reaches a specified $x$- or $y$-position.
&
\textbf{EK 4.1.A.1}: A parametric function $f(t) = (x(t),\, y(t))$ in
$\mathbb{R}^2$ consists of two equations with dependent variables $x$ and
$y$, each expressed as a function of a single independent variable $t$
called the \emph{parameter}. \\[4pt]

\textbf{2.B} \textit{Construct equivalent representations} ---
  Build a numerical table of $(x_i, y_i)$ pairs and connect the plotted
  points, in order of increasing $t$, to produce the graph of a
  parametric function.
&
\textbf{EK 4.1.A.2}: The coordinates $(x_i,\, y_i)$ at time $t_i$ can
be expressed as the single parametric function
$f(t) = (x(t),\, y(t))$. \\[4pt]

&
\textbf{EK 4.1.A.3}: A numerical table is generated by evaluating $x_i$
and $y_i$ at several values of $t_i$ within the domain. \\[4pt]

&
\textbf{EK 4.1.A.4}: A graph is constructed by connecting the table
points in order of increasing $t$. \\[4pt]

&
\textbf{EK 4.1.A.5}: A restricted domain results in start and end
points on the parametric curve. \\
\end{tabularx}
\end{skillbox}

\vspace{0.08in}

%----------------------------------------------------------------------
% Vocabulary
%----------------------------------------------------------------------
\begin{skillbox}[{Vocabulary, Concepts, \& Theorems}]{greenbox}
\begin{tabularx}{\linewidth}{@{} p{0.24\linewidth} X @{}}
\textbf{Term} & \textbf{Definition / Note} \\
\midrule
parameter &
  The independent variable $t$ in a parametric function; it controls
  both $x$ and $y$ simultaneously. \\[4pt]
parametric function &
  A function $f(t) = (x(t),\, y(t))$ in $\mathbb{R}^2$ consisting of
  two component equations, both in terms of the parameter $t$. \\[4pt]
parametric curve &
  The set of ordered pairs $(x(t),\, y(t))$ traced in the $xy$-plane
  as $t$ increases through its domain; the curve has an
  \emph{orientation} (direction) determined by increasing $t$. \\[4pt]
start / end point &
  When the domain of $t$ is restricted to $[a, b]$, the curve begins
  at $(x(a),\, y(a))$ and ends at $(x(b),\, y(b))$. \\
\end{tabularx}
\end{skillbox}

\vspace{0.08in}

%----------------------------------------------------------------------
% Activate Prior Knowledge
%----------------------------------------------------------------------
\begin{skillbox}[Activate Prior Knowledge \& Spiral Review]{sky}
\begin{tabularx}{\linewidth}{@{}X c@{}}
\begin{minipage}[t]{0.96\linewidth}\vspace{0pt}
\textbf{Spiral topics activated during Warm-Up (first 8 minutes):}
\begin{itemize}[leftmargin=1.4em, itemsep=2pt]
  \item \textbf{Evaluating algebraic expressions}: Substitute values of
    $t$ into polynomial expressions --- the core computation in building
    a parametric table.
  \item \textbf{Paired function evaluation}: Evaluate two functions at
    the same input and record both outputs together as an ordered pair
    --- directly previews the parametric table structure.
  \item \textbf{Ordered pairs}: Recall that a point in the plane is
    described by an ordered pair $(x, y)$; parametric functions output
    exactly such a pair for each $t$.
\end{itemize}
\end{minipage}
&
\end{tabularx}
\end{skillbox}

\vspace{0.08in}

%----------------------------------------------------------------------
% Hook
%----------------------------------------------------------------------
\begin{skillbox}[Hook --- Entry Event]{sky}
\textbf{Scenario:} A drone launches from a pad and its eastward position
(meters east of the pad) at time $t$ seconds is $x(t) = t^2 - 4$, and
its northward position is $y(t) = 2t$ for $0 \le t \le 4$.

\textbf{Discussion prompts:}
\begin{enumerate}[leftmargin=1.4em, itemsep=3pt]
  \item \textit{``Can you write a single equation $y = f(x)$ that
    describes the drone's path? What would you lose if you did?''}
  \item \textit{``What information does the parameter $t$ carry that a
    plain $xy$-equation does not?''}
  \item \textit{``Where does the drone's path start and where does it
    end?''}
\end{enumerate}

\textbf{Key tension to surface:} A single $y = f(x)$ equation loses
track of \emph{time} and \emph{direction}. The parametric form tells us
not just \emph{where} the drone goes, but \emph{when} --- and gives the
curve an orientation.

\textbf{Transition:} ``Today we build the table and graph from two
component equations, and we will see how the single parameter $t$
controls both coordinates at once.''
\end{skillbox}

\vspace{0.08in}

%----------------------------------------------------------------------
% Lesson
%----------------------------------------------------------------------
\begin{skillbox}[Lesson --- Instruction Sequence]{sky}
\begin{multicols}{2}

\textbf{Step 1 --- Define Parametric Functions} (8 min)

Using the drone hook, formally write $f(t) = (x(t),\, y(t))$ on the
board and highlight that $t$ is the \emph{independent} variable even
though it does not appear on either axis of the graph. Contrast with
$y = g(x)$: parametric functions output a \emph{pair} of values. Cover
EK 4.1.A.1--A.2. Ask: ``Is $f(t) = (x(t),\, y(t))$ a function? What
is its input? What is its output?''

\columnbreak

\textbf{Step 2 --- Build a Table of Values} (12 min)

Model with $f(t) = (t+1,\, t^2 - 3)$, $-2 \le t \le 2$. Compute
$(x, y)$ at $t = -2, -1, 0, 1, 2$ step by step, filling in the table
column-by-column (all $x$ values first, then all $y$ values, then the
ordered pairs). EK 4.1.A.3. Students practice with the second function
in guided notes.

\end{multicols}

\begin{multicols}{2}

\textbf{Step 3 --- Graph the Curve} (10 min)

Plot the five ordered pairs from Step 2 on the pre-printed grid in the
guided notes. Connect the points \emph{in order of increasing $t$},
drawing arrows to show orientation. Identify the start point $(-1, 1)$
at $t = -2$ and end point $(3, 1)$ at $t = 2$. EK 4.1.A.4--A.5.
Emphasize: order matters; the curve has \emph{direction}.

\columnbreak

\textbf{Step 4 --- Solving for $t$ (Skill 1.A)} (10 min)

Pose: ``The curve crosses the $y$-axis when $x = 0$. When does that
happen?'' Model solving $x(t) = t+1 = 0 \Rightarrow t = -1$, then
evaluating $y(-1) = (-1)^2 - 3 = -2$. The curve crosses the $y$-axis
at $(0, -2)$. Students solve: find all $t$ in the domain where $y(t) = -2$.

\end{multicols}

\smallskip
\textbf{Pacing note:} Steps 1--2 in Period 1; Steps 3--4 + activity in Period 2.
\end{skillbox}

\vspace{0.08in}

%----------------------------------------------------------------------
% Active Monitoring
%----------------------------------------------------------------------
\begin{skillbox}[Active Monitoring]{redbox}
\begin{itemize}[leftmargin=1.4em, itemsep=3pt]
  \item \textbf{Column confusion}: Students may mix up which component
    gives $x$ and which gives $y$. Prompt: \textit{``Which equation are
    you evaluating right now? What axis does that output go on?''}
  \item \textbf{Ordering errors}: Students may plot points without
    regard to the order of $t$. Prompt: \textit{``Trace the curve by
    connecting the points in the order you computed them. Where do you
    start?''}
  \item \textbf{Domain restriction}: Students may draw the curve beyond
    its endpoints. Ask: \textit{``Does your graph go on forever, or
    does it have endpoints? What are they?''}
  \item \textbf{Solving for $t$}: Students solve $x(t) = c$ and report
    $t$ as the final answer. Prompt: \textit{``The question asked for a
    \emph{position} (an ordered pair), not a time. What are the
    coordinates at that moment?''}
\end{itemize}
\end{skillbox}

\vspace{0.08in}

%----------------------------------------------------------------------
% Group Work & Differentiation
%----------------------------------------------------------------------
\begin{skillbox}[Group Work \& Differentiation]{redbox}
Students work in groups of 3--4 on the tiered activity (assign by
teacher or allow student self-selection with guidance).

\begin{itemize}[leftmargin=1.4em, itemsep=4pt]
  \item \textbf{Tier R (Remediate)}: Given a partially completed table
    for $f(t) = (t^2,\, t-1)$ on $-2 \le t \le 3$, students fill in
    the missing values and identify the start and end points without
    graphing.
  \item \textbf{Tier A (Accelerate)}: Students complete the full table,
    identify start and end points, and solve for $t$ when the curve
    crosses the $x$-axis and when it crosses the $y$-axis.
  \item \textbf{Tier E (Extend)}: Students observe that $t = -1$ and
    $t = 1$ both give $x(t) = 1$ but different $y$-values, and explain
    what this means for the curve (it passes through $x = 1$ twice)
    and why the curve cannot be expressed as $y = g(x)$ on the full
    domain.
\end{itemize}
\end{skillbox}

\vspace{0.08in}

%----------------------------------------------------------------------
% Individual Work & Assessment
%----------------------------------------------------------------------
\begin{skillbox}[Individual Work \& Assessment]{redbox}
Exit ticket administered independently (final 5 minutes of Period 2).

\begin{enumerate}[leftmargin=1.4em, itemsep=4pt]
  \item Given $f(t) = (2t-1,\; 3-t)$ for $0 \le t \le 3$, complete a
    four-row table of values and identify the start and end points.
  \item Solve: for what value of $t$ does the curve cross the $y$-axis?
    State the coordinates of that point.
\end{enumerate}

\textbf{Data use:} Students who cannot evaluate both components at a
single $t$-value need more practice with function evaluation; address
in the opening minutes of Lesson 4.2.
\end{skillbox}

\vspace{0.08in}

%----------------------------------------------------------------------
% Reinforcement & Extension
%----------------------------------------------------------------------
\begin{skillbox}[Reinforcement \& Extension]{goldbox}
\textbf{Homework (Lesson 8.1):} 5--6 problems covering: building tables
from two-component parametric equations, identifying start and end
points from a restricted domain, solving for $t$ given an $x$- or
$y$-value, and 1 AP-style multiple-choice item.

\textbf{Preview of Lesson 4.2 --- Parametric Functions Modeling Planar
Motion:} Ask students: \textit{``If $x(t)$ and $y(t)$ each have a
maximum and minimum value, how could you find when a moving particle
reaches its leftmost or highest point?''} Lesson 4.2 applies today's
table-and-solve skills to identify horizontal and vertical extrema of
particle motion in the plane.
\end{skillbox}

\end{document}
